\documentclass[11pt]{article}
% The preamble every project's paper shares: % The preamble every project's paper shares: % The preamble every project's paper shares: \input{../../preamble} from
% projects/<project>/paper/. Boards come from chessboard and xskak, plots from
% pgfplots; every number a paper prints comes from its generated.tex.
\usepackage[LSBC4,T1]{fontenc}
\usepackage[utf8]{inputenc}
\usepackage{lmodern}
\usepackage[margin=1in]{geometry}
\usepackage{microtype}
\usepackage{amsmath}
\usepackage{booktabs}
\usepackage{tabularx}
\usepackage{array}
\usepackage{xcolor}
\usepackage[obeyspaces,spaces]{url}
\usepackage{xskak}
\usepackage{pgfplots}
\pgfplotsset{compat=1.18}
\usepackage{caption}
\usepackage{subcaption}
\usepackage{enumitem}
\usepackage[section]{placeins}
\usepackage[hidelinks]{hyperref}

\setchessboard{boardfontencoding=LSBC4, boardfontfamily=skaknew, blackfieldcolor=black!22,
  whitefieldcolor=white, setfontcolors, showmover=false, label=true, labelfontsize=6pt,
  boardfontsize=14pt}

% A move or line in UCI, and a position in FEN, both breakable at a slash.
\newcommand{\uci}[1]{\texttt{#1}}
\newcommand{\fen}[1]{{\footnotesize\path{#1}}}
\newcolumntype{L}{>{\raggedright\arraybackslash}X}

\setlength{\parskip}{0.4em}
\setlength{\parindent}{0pt}
 from
% projects/<project>/paper/. Boards come from chessboard and xskak, plots from
% pgfplots; every number a paper prints comes from its generated.tex.
\usepackage[LSBC4,T1]{fontenc}
\usepackage[utf8]{inputenc}
\usepackage{lmodern}
\usepackage[margin=1in]{geometry}
\usepackage{microtype}
\usepackage{amsmath}
\usepackage{booktabs}
\usepackage{tabularx}
\usepackage{array}
\usepackage{xcolor}
\usepackage[obeyspaces,spaces]{url}
\usepackage{xskak}
\usepackage{pgfplots}
\pgfplotsset{compat=1.18}
\usepackage{caption}
\usepackage{subcaption}
\usepackage{enumitem}
\usepackage[section]{placeins}
\usepackage[hidelinks]{hyperref}

\setchessboard{boardfontencoding=LSBC4, boardfontfamily=skaknew, blackfieldcolor=black!22,
  whitefieldcolor=white, setfontcolors, showmover=false, label=true, labelfontsize=6pt,
  boardfontsize=14pt}

% A move or line in UCI, and a position in FEN, both breakable at a slash.
\newcommand{\uci}[1]{\texttt{#1}}
\newcommand{\fen}[1]{{\footnotesize\path{#1}}}
\newcolumntype{L}{>{\raggedright\arraybackslash}X}

\setlength{\parskip}{0.4em}
\setlength{\parindent}{0pt}
 from
% projects/<project>/paper/. Boards come from chessboard and xskak, plots from
% pgfplots; every number a paper prints comes from its generated.tex.
\usepackage[LSBC4,T1]{fontenc}
\usepackage[utf8]{inputenc}
\usepackage{lmodern}
\usepackage[margin=1in]{geometry}
\usepackage{microtype}
\usepackage{amsmath}
\usepackage{booktabs}
\usepackage{tabularx}
\usepackage{array}
\usepackage{xcolor}
\usepackage[obeyspaces,spaces]{url}
\usepackage{xskak}
\usepackage{pgfplots}
\pgfplotsset{compat=1.18}
\usepackage{caption}
\usepackage{subcaption}
\usepackage{enumitem}
\usepackage[section]{placeins}
\usepackage[hidelinks]{hyperref}

\setchessboard{boardfontencoding=LSBC4, boardfontfamily=skaknew, blackfieldcolor=black!22,
  whitefieldcolor=white, setfontcolors, showmover=false, label=true, labelfontsize=6pt,
  boardfontsize=14pt}

% A move or line in UCI, and a position in FEN, both breakable at a slash.
\newcommand{\uci}[1]{\texttt{#1}}
\newcommand{\fen}[1]{{\footnotesize\path{#1}}}
\newcolumntype{L}{>{\raggedright\arraybackslash}X}

\setlength{\parskip}{0.4em}
\setlength{\parindent}{0pt}

% Written by projects/book-advice/paper/generate.mjs; do not edit by hand.

\newcommand{\AppCommit}{50de20ccbf}
\newcommand{\Month}{September 2026}
\newcommand{\PlayerGamesRead}{148,000}
\newcommand{\GamesRead}{74,000}
\newcommand{\ArchiveGames}{639}
\newcommand{\HisPlayerGames}{39,903}
\newcommand{\AbovePlayerGames}{34,973}
\newcommand{\HisLo}{1250}
\newcommand{\HisHi}{1449}
\newcommand{\AboveLo}{1450}
\newcommand{\AboveHi}{1649}
\newcommand{\CheckNodes}{120,000}
\newcommand{\BestNodes}{600,000}
\newcommand{\VerifyNodes}{750,000}
\newcommand{\GradeNodes}{250,000}
\newcommand{\PlanNodes}{600,000}
\newcommand{\GiftFloor}{20}
\newcommand{\HighWp}{75}
\newcommand{\SwingPlies}{six}
\newcommand{\SwingWp}{30}
\newcommand{\OneAHisEvents}{144,404}
\newcommand{\OneAHisPunished}{62.0}
\newcommand{\OneAHisTrap}{6.4}
\newcommand{\OneAHisMissed}{31.1}
\newcommand{\OneAHisTrapOfFailures}{17.1\% {\footnotesize (16.8--17.4)}}
\newcommand{\OneAHisMissedOfFailures}{82.9\% {\footnotesize (82.5--83.2)}}
\newcommand{\OneAHisGiftTrapOfFailures}{21.8\% {\footnotesize (21.3--22.4)}}
\newcommand{\OneAHisGiftPunished}{65.9}
\newcommand{\OneAHisNextSerious}{21.7\% {\footnotesize (21.5--22.0)}}
\newcommand{\OneAHisHighOtherSerious}{9.6\% {\footnotesize (9.5--9.7)}}
\newcommand{\OneAAboveEvents}{121,595}
\newcommand{\OneAAbovePunished}{63.4}
\newcommand{\OneAAboveTrap}{6.0}
\newcommand{\OneAAboveMissed}{30.2}
\newcommand{\OneAAboveTrapOfFailures}{16.7\% {\footnotesize (16.3--17.0)}}
\newcommand{\OneAAboveMissedOfFailures}{83.3\% {\footnotesize (83.0--83.7)}}
\newcommand{\OneAAboveGiftTrapOfFailures}{21.5\% {\footnotesize (20.9--22.1)}}
\newcommand{\OneAAboveGiftPunished}{67.0}
\newcommand{\OneAAboveNextSerious}{20.5\% {\footnotesize (20.2--20.8)}}
\newcommand{\OneAAboveHighOtherSerious}{9.1\% {\footnotesize (9.0--9.2)}}
\newcommand{\OneAArchiveEvents}{1,927}
\newcommand{\OneAArchivePunished}{65.3}
\newcommand{\OneAArchiveTrap}{6.4}
\newcommand{\OneAArchiveMissed}{27.6}
\newcommand{\OneAArchiveTrapOfFailures}{18.9\% {\footnotesize (16.1--22.1)}}
\newcommand{\OneAArchiveMissedOfFailures}{81.1\% {\footnotesize (77.9--83.9)}}
\newcommand{\OneAArchiveGiftTrapOfFailures}{24.6\% {\footnotesize (19.3--30.9)}}
\newcommand{\OneAArchiveGiftPunished}{72.8}
\newcommand{\OneAArchiveNextSerious}{18.0\% {\footnotesize (15.9--20.3)}}
\newcommand{\OneAArchiveHighOtherSerious}{8.9\% {\footnotesize (8.0--9.9)}}
\newcommand{\OneATable}{%
\begin{tabular}{lrrrrll}
\toprule
 & & \multicolumn{3}{c}{share of the replies} & \multicolumn{2}{c}{share of the failures} \\
\cmidrule(lr){3-5}\cmidrule(lr){6-7}
players & replies & punished & a trap taken & missed & a trap taken & missed \\
\midrule
his band & 144,404 & 62.0\% & 6.4\% & 31.1\% & 17.1\% {\footnotesize (16.8--17.4)} & 82.9\% {\footnotesize (82.5--83.2)} \\
the band above & 121,595 & 63.4\% & 6.0\% & 30.2\% & 16.7\% {\footnotesize (16.3--17.0)} & 83.3\% {\footnotesize (83.0--83.7)} \\
his games & 1,927 & 65.3\% & 6.4\% & 27.6\% & 18.9\% {\footnotesize (16.1--22.1)} & 81.1\% {\footnotesize (77.9--83.9)} \\
\bottomrule
\end{tabular}
}
\newcommand{\OneASampleMisses}{2,000}
\newcommand{\OneASampleOffers}{3,000}
\newcommand{\OneAHisMissesSearched}{2,000}
\newcommand{\OneAHisKindHangingPiece}{10.3\% {\footnotesize (9.1--11.8)}}
\newcommand{\OneAHisKindMate}{0.8\% {\footnotesize (0.5--1.3)}}
\newcommand{\OneAHisKindFork}{1.6\% {\footnotesize (1.1--2.2)}}
\newcommand{\OneAHisKindPinSkewerDiscovered}{5.7\% {\footnotesize (4.8--6.8)}}
\newcommand{\OneAHisKindStrike}{0.2\% {\footnotesize (0.1--0.5)}}
\newcommand{\OneAHisKindDefensiveMove}{19.6\% {\footnotesize (17.9--21.3)}}
\newcommand{\OneAHisKindUnnamed}{61.9\% {\footnotesize (59.7--63.9)}}
\newcommand{\OneAHisOfferedSearched}{3,000}
\newcommand{\OneAHisOfferedTraps}{12.6\% {\footnotesize (11.5--13.8)}}
\newcommand{\OneAHisTookTrap}{27.8\% {\footnotesize (23.5--32.5)}}
\newcommand{\OneAHisLeftGift}{34.1\% {\footnotesize (32.3--35.9)}}
\newcommand{\OneAHisOfferedTrapKnight}{17.5\% {\footnotesize (15.3--20.0)}}
\newcommand{\OneAHisOfferedTrapBishop}{13.3\% {\footnotesize (11.1--15.9)}}
\newcommand{\OneAHisOfferedTrapRook}{10.0\% {\footnotesize (8.1--12.2)}}
\newcommand{\OneAHisOfferedTrapQueen}{3.2\% {\footnotesize (1.9--5.6)}}
\newcommand{\OneAAboveMissesSearched}{2,000}
\newcommand{\OneAAboveKindHangingPiece}{9.5\% {\footnotesize (8.3--10.9)}}
\newcommand{\OneAAboveKindMate}{0.6\% {\footnotesize (0.3--1.1)}}
\newcommand{\OneAAboveKindFork}{1.7\% {\footnotesize (1.2--2.4)}}
\newcommand{\OneAAboveKindPinSkewerDiscovered}{5.0\% {\footnotesize (4.1--6.0)}}
\newcommand{\OneAAboveKindStrike}{0.1\% {\footnotesize (0.0--0.4)}}
\newcommand{\OneAAboveKindDefensiveMove}{16.4\% {\footnotesize (14.9--18.1)}}
\newcommand{\OneAAboveKindUnnamed}{66.7\% {\footnotesize (64.6--68.7)}}
\newcommand{\OneAAboveOfferedSearched}{3,000}
\newcommand{\OneAAboveOfferedTraps}{14.3\% {\footnotesize (13.1--15.6)}}
\newcommand{\OneAAboveTookTrap}{23.8\% {\footnotesize (20.0--28.1)}}
\newcommand{\OneAAboveLeftGift}{32.7\% {\footnotesize (30.9--34.5)}}
\newcommand{\OneAAboveOfferedTrapKnight}{19.1\% {\footnotesize (16.8--21.6)}}
\newcommand{\OneAAboveOfferedTrapBishop}{16.9\% {\footnotesize (14.3--19.8)}}
\newcommand{\OneAAboveOfferedTrapRook}{11.0\% {\footnotesize (9.1--13.3)}}
\newcommand{\OneAAboveOfferedTrapQueen}{4.9\% {\footnotesize (3.2--7.3)}}
\newcommand{\OneAArchiveMissesSearched}{531}
\newcommand{\OneAArchiveKindHangingPiece}{10.9\% {\footnotesize (8.6--13.9)}}
\newcommand{\OneAArchiveKindMate}{1.5\% {\footnotesize (0.8--2.9)}}
\newcommand{\OneAArchiveKindFork}{2.5\% {\footnotesize (1.4--4.1)}}
\newcommand{\OneAArchiveKindPinSkewerDiscovered}{3.8\% {\footnotesize (2.5--5.8)}}
\newcommand{\OneAArchiveKindStrike}{0.0\% {\footnotesize (0.0--0.7)}}
\newcommand{\OneAArchiveKindDefensiveMove}{19.4\% {\footnotesize (16.3--23.0)}}
\newcommand{\OneAArchiveKindUnnamed}{62.0\% {\footnotesize (57.8--66.0)}}
\newcommand{\OneAArchiveOfferedSearched}{228}
\newcommand{\OneAArchiveOfferedTraps}{13.6\% {\footnotesize (9.8--18.6)}}
\newcommand{\OneAArchiveTookTrap}{29.0\% {\footnotesize (16.1--46.6)}}
\newcommand{\OneAArchiveLeftGift}{29.4\% {\footnotesize (23.5--36.1)}}
\newcommand{\OneAArchiveOfferedTrapKnight}{15.1\% {\footnotesize (9.0--24.2)}}
\newcommand{\OneAArchiveOfferedTrapBishop}{14.0\% {\footnotesize (6.6--27.3)}}
\newcommand{\OneAArchiveOfferedTrapRook}{9.0\% {\footnotesize (4.2--18.2)}}
\newcommand{\OneAArchiveOfferedTrapQueen}{18.8\% {\footnotesize (8.9--35.3)}}
\newcommand{\OneAHisTrapShare}{17}
\newcommand{\OneAHisPunishedShare}{62}
\newcommand{\OneAKindsTable}{%
\begin{tabular}{llll}
\toprule
the engine's best reply & his band & the band above & his games \\
\midrule
misses searched & 2,000 & 2,000 & 531 \\
\midrule
the hanging piece & 10.3\% {\footnotesize (9.1--11.8)} & 9.5\% {\footnotesize (8.3--10.9)} & 10.9\% {\footnotesize (8.6--13.9)} \\
a mate & 0.8\% {\footnotesize (0.5--1.3)} & 0.6\% {\footnotesize (0.3--1.1)} & 1.5\% {\footnotesize (0.8--2.9)} \\
a fork & 1.6\% {\footnotesize (1.1--2.2)} & 1.7\% {\footnotesize (1.2--2.4)} & 2.5\% {\footnotesize (1.4--4.1)} \\
a pin, skewer or discovered attack & 5.7\% {\footnotesize (4.8--6.8)} & 5.0\% {\footnotesize (4.1--6.0)} & 3.8\% {\footnotesize (2.5--5.8)} \\
the f2 or f7 strike & 0.2\% {\footnotesize (0.1--0.5)} & 0.1\% {\footnotesize (0.0--0.4)} & 0.0\% {\footnotesize (0.0--0.7)} \\
saving a loose piece of his own & 19.6\% {\footnotesize (17.9--21.3)} & 16.4\% {\footnotesize (14.9--18.1)} & 19.4\% {\footnotesize (16.3--23.0)} \\
a move it does not name & 61.9\% {\footnotesize (59.7--63.9)} & 66.7\% {\footnotesize (64.6--68.7)} & 62.0\% {\footnotesize (57.8--66.0)} \\
\bottomrule
\end{tabular}
}
\newcommand{\KingShowsBar}{2}
\newcommand{\KingLiftBar}{1.5}
\newcommand{\KingFloorBar}{1.25}
\newcommand{\KingHisPositions}{8,000}
\newcommand{\KingHisBase}{12.4\% {\footnotesize (11.7--13.2)}}
\newcommand{\KingAbovePositions}{8,000}
\newcommand{\KingAboveBase}{10.8\% {\footnotesize (10.2--11.6)}}
\newcommand{\KingArchivePositions}{6,215}
\newcommand{\KingArchiveBase}{9.3\% {\footnotesize (8.6--10.0)}}
\newcommand{\KingHisMateBase}{2.3}
\newcommand{\KingHisSwingBase}{10.3}
\newcommand{\KingTable}{%
\begin{tabular}{lrlrcrr}
\toprule
 & \multicolumn{4}{c}{his band} & above & his games \\
\cmidrule(lr){2-5}
sign & shows & dangerous when shown & lift & clears & lift & lift \\
\midrule
every position & & 12.4\% {\footnotesize (11.7--13.2)} & & & & \\
\midrule
no safe squares & 3.6\% & 20.7\% {\footnotesize (16.4--25.7)} & 1.66 & yes & 1.80 & 2.55 \\
queen on the exits & 6.2\% & 30.3\% {\footnotesize (26.4--34.5)} & 2.43 & yes & 2.43 & 3.15 \\
bishop aimed in & 47.3\% & 14.4\% {\footnotesize (13.4--15.6)} & 1.16 & no & 1.19 & 1.01 \\
approaching knight & 72.2\% & 12.1\% {\footnotesize (11.3--12.9)} & 0.97 & no & 0.94 & 0.96 \\
rook lift & 24.1\% & 18.8\% {\footnotesize (17.1--20.6)} & 1.51 & yes & 1.58 & 1.72 \\
hook pawn & 57.8\% & 10.8\% {\footnotesize (9.9--11.7)} & 0.87 & no & 0.83 & 0.86 \\
opposite-side storm & 3.6\% & 19.3\% {\footnotesize (15.1--24.3)} & 1.55 & no & 1.43 & 1.85 \\
plus two & 6.5\% & 21.1\% {\footnotesize (17.9--24.9)} & 1.70 & yes & 1.60 & 2.19 \\
\bottomrule
\end{tabular}
}
\newcommand{\KingCleared}{four}
\newcommand{\KingClearedCap}{Four}
\newcommand{\KingStormLift}{1.55}
\newcommand{\KingStormLow}{15.1}
\newcommand{\KingFloor}{15.6}
\newcommand{\KingNoSafeMateLift}{5.1}
\newcommand{\KingNoSafeSwingLift}{0.93}
\newcommand{\BookPlies}{12}
\newcommand{\KingEarlyPositions}{1,755}
\newcommand{\KingEarlyBase}{6.0}
\newcommand{\KingEarlyAny}{7.6\% {\footnotesize (4.2--13.4)}}
\newcommand{\KingEarlyNoSafeShows}{6.2}
\newcommand{\KingLatePositions}{6,245}
\newcommand{\KingLateBase}{14.3\% {\footnotesize (13.4--15.2)}}
\newcommand{\KingLateTable}{%
\begin{tabular}{lrlrcrrr}
\toprule
 & \multicolumn{4}{c}{his band, after ply 12} & \multicolumn{3}{c}{against the ply alone} \\
\cmidrule(lr){2-5}\cmidrule(lr){6-8}
sign & shows & dangerous when shown & lift & clears & his band & above & his games \\
\midrule
every position & & 14.3\% {\footnotesize (13.4--15.2)} & & & & & \\
\midrule
no safe squares & 2.9\% & 29.1\% {\footnotesize (23.0--36.1)} & 2.04 & yes & 1.88 & 2.01 & 2.71 \\
queen on the exits & 7.6\% & 31.1\% {\footnotesize (27.1--35.4)} & 2.18 & yes & 1.95 & 1.79 & 2.26 \\
rook lift & 30.7\% & 18.8\% {\footnotesize (17.1--20.6)} & 1.31 & no & 1.17 & 1.14 & 1.16 \\
plus two & 8.3\% & 21.3\% {\footnotesize (18.0--25.1)} & 1.49 & no & 1.32 & 1.17 & 1.44 \\
\bottomrule
\end{tabular}
}
\newcommand{\KingBinFirst}{11.1}
\newcommand{\KingBinLast}{18.7}
\newcommand{\KingBins}{13--20, 21--30, 31--40 and 41 and later}
\newcommand{\KingBinFirstPlies}{13 to 20}
\newcommand{\KingBinLastFrom}{41}
\newcommand{\KingLateNoSafeSquares}{2.04}
\newcommand{\KingLateQueenOnExits}{2.18}
\newcommand{\KingLateRookLift}{1.31}
\newcommand{\KingLatePlusTwo}{1.49}
\newcommand{\KingNotesPerGame}{0.46}
\newcommand{\KingNotesSilent}{65}
\newcommand{\KingNotesHisBand}{0.53}
\newcommand{\KingNotesAllPerGame}{1.61}
\newcommand{\KingOnsetDanger}{26.7\% {\footnotesize (20.5--33.9)}}
\newcommand{\KingNoneDanger}{9.7\% {\footnotesize (8.8--10.6)}}
\newcommand{\KingAllOnsetDanger}{16.4\% {\footnotesize (13.4--19.9)}}
\newcommand{\TradeShareBar}{5}
\newcommand{\TradeFaultBar}{2}
\newcommand{\TradeMobilityGap}{three}
\newcommand{\TradeHisTrades}{1,500}
\newcommand{\TradeHisPraised}{73}
\newcommand{\TradeHisFaulted}{2.8\% {\footnotesize (2.0--4.0)}}
\newcommand{\TradeHisFaultedCount}{31 of 1,094}
\newcommand{\TradeAboveTrades}{1,500}
\newcommand{\TradeAbovePraised}{73}
\newcommand{\TradeAboveFaulted}{2.8\% {\footnotesize (2.0--4.0)}}
\newcommand{\TradeAboveFaultedCount}{31 of 1,096}
\newcommand{\TradeArchiveTrades}{89}
\newcommand{\TradeArchivePraised}{79}
\newcommand{\TradeArchiveFaulted}{5.7\% {\footnotesize (2.2--13.8)}}
\newcommand{\TradeArchiveFaultedCount}{4 of 70}
\newcommand{\TradeTable}{%
\begin{tabular}{lrllrll}
\toprule
 & \multicolumn{3}{c}{his band} & \multicolumn{3}{c}{his games} \\
\cmidrule(lr){2-4}\cmidrule(lr){5-7}
case & share & faulted & others faulted & share & faulted & others faulted \\
\midrule
active for passive & 28.1\% & 1.6\% & 3.3\% & 32.9\% & 1 of 23 & 3 of 47 \\
an attacker traded & 2.8\% & 0.0\% & 2.9\% & 4.3\% & 0 of 3 & 4 of 67 \\
queens off while attacking & 3.3\% & 11.1\% & 2.5\% & 2.9\% & 0 of 2 & 4 of 68 \\
a worse ending & 1.3\% & all & 1.6\% & 4.3\% & all & 1 of 67 \\
\bottomrule
\end{tabular}
}
\newcommand{\TradeActiveShare}{28}
\newcommand{\TradeActiveFaulted}{1.6}
\newcommand{\TradeActiveOthers}{3.3}
\newcommand{\TradeQueensFaulted}{11.1}
\newcommand{\TradeQueensShare}{3.3}
\newcommand{\TradeLineFires}{32}
\newcommand{\TradeLineFaultedAhead}{840}
\newcommand{\TradeLineActive}{7 of 32}
\newcommand{\TradeLineAttacker}{2 of 32}
\newcommand{\TradeLineQueens}{2 of 32}
\newcommand{\RetreatsDrawGames}{3,000}
\newcommand{\RetreatsDrawArchive}{6,000}
\newcommand{\RetreatsMustFind}{10}
\newcommand{\RetreatsBins}{10 to 20, 20 to 40 and 40 or more}
\newcommand{\RetreatsHisPositions}{29,887}
\newcommand{\RetreatsHisMustFindPositions}{3,101}
\newcommand{\RetreatsHisBackward}{36.5\% {\footnotesize (30.7--42.7)}}
\newcommand{\RetreatsHisForward}{46.7\% {\footnotesize (43.0--50.3)}}
\newcommand{\RetreatsHisDiff}{$-$10.2 {\footnotesize ($-$17.0 to $-$3.0)}}
\newcommand{\RetreatsHisKnightBack}{32.8\% {\footnotesize (22.1--45.6)}}
\newcommand{\RetreatsHisKnightFwd}{45.0\% {\footnotesize (36.6--53.6)}}
\newcommand{\RetreatsHisAllBackward}{30.8\% {\footnotesize (26.4--35.6)}}
\newcommand{\RetreatsHisAllForward}{27.9\% {\footnotesize (26.1--29.7)}}
\newcommand{\RetreatsHisMaiaDiff}{6.8}
\newcommand{\RetreatsHisMaiaCi}{2.4 to 11.3}
\newcommand{\RetreatsAbovePositions}{29,904}
\newcommand{\RetreatsAboveMustFindPositions}{3,050}
\newcommand{\RetreatsAboveBackward}{36.3\% {\footnotesize (30.7--42.4)}}
\newcommand{\RetreatsAboveForward}{47.2\% {\footnotesize (43.4--51.0)}}
\newcommand{\RetreatsAboveDiff}{$-$10.8 {\footnotesize ($-$17.6 to $-$3.7)}}
\newcommand{\RetreatsAboveKnightBack}{40.0\% {\footnotesize (28.6--52.6)}}
\newcommand{\RetreatsAboveKnightFwd}{42.0\% {\footnotesize (34.2--50.1)}}
\newcommand{\RetreatsAboveAllBackward}{30.7\% {\footnotesize (26.4--35.4)}}
\newcommand{\RetreatsAboveAllForward}{27.1\% {\footnotesize (25.3--28.9)}}
\newcommand{\RetreatsAboveMaiaDiff}{10.0}
\newcommand{\RetreatsAboveMaiaCi}{5.3 to 14.6}
\newcommand{\RetreatsArchivePositions}{6,000}
\newcommand{\RetreatsArchiveMustFindPositions}{607}
\newcommand{\RetreatsArchiveBackward}{40.0\% {\footnotesize (27.6--53.8)}}
\newcommand{\RetreatsArchiveForward}{40.2\% {\footnotesize (31.9--49.0)}}
\newcommand{\RetreatsArchiveDiff}{$-$0.2 {\footnotesize ($-$15.4 to 15.9)}}
\newcommand{\RetreatsArchiveKnightBack}{46.2\% {\footnotesize (23.2--70.9)}}
\newcommand{\RetreatsArchiveKnightFwd}{23.8\% {\footnotesize (10.6--45.1)}}
\newcommand{\RetreatsArchiveAllBackward}{31.1\% {\footnotesize (22.5--41.3)}}
\newcommand{\RetreatsArchiveAllForward}{25.7\% {\footnotesize (21.7--30.1)}}
\newcommand{\RetreatsArchiveMaiaDiff}{8.8}
\newcommand{\RetreatsArchiveMaiaCi}{$-$1.8 to 19.6}
\newcommand{\RetreatsTable}{%
\setlength{\tabcolsep}{4pt}
\begin{tabular}{lrllll}
\toprule
 & & \multicolumn{3}{c}{quiet best moves missed} & \\
\cmidrule(lr){3-5}
players & must-find & backward & sideways & forward & backward $-$ forward \\
\midrule
his band & 3,101 & 36.5\% {\footnotesize (30.7--42.7)} & 37.8\% {\footnotesize (31.4--44.7)} & 46.7\% {\footnotesize (43.0--50.3)} & $-$10.2 {\footnotesize ($-$17.0 to $-$3.0)} \\
the band above & 3,050 & 36.3\% {\footnotesize (30.7--42.4)} & 40.4\% {\footnotesize (34.0--47.2)} & 47.2\% {\footnotesize (43.4--51.0)} & $-$10.8 {\footnotesize ($-$17.6 to $-$3.7)} \\
his games & 607 & 40.0\% {\footnotesize (27.6--53.8)} & 35.4\% {\footnotesize (23.4--49.6)} & 40.2\% {\footnotesize (31.9--49.0)} & $-$0.2 {\footnotesize ($-$15.4 to 15.9)} \\
\bottomrule
\end{tabular}
}
\newcommand{\RetreatsMaiaTable}{%
\begin{tabular}{lrrrr}
\toprule
 & \multicolumn{2}{c}{backward} & \multicolumn{2}{c}{forward} \\
\cmidrule(lr){2-3}\cmidrule(lr){4-5}
best move's margin (win\%) & positions & Maia-2 & positions & Maia-2 \\
\midrule
10--20 & 147 & 51.1\% & 466 & 43.8\% \\
20--40 & 88 & 64.4\% & 206 & 57.0\% \\
40 or more & 9 & 61.5\% & 44 & 70.9\% \\
\midrule
weighted by the backward moves & & \multicolumn{3}{l}{backward $-$ forward: 6.8 points (2.4 to 11.3)} \\
\bottomrule
\end{tabular}
}
\newcommand{\RetreatsLastBinBackward}{nine}
\newcommand{\RetreatsLastBinFrom}{40}
\newcommand{\RetreatsFirstBinsBackward}{235}
\newcommand{\RetreatsHisQuietBackward}{244}
\newcommand{\RetreatsFwdQuietShare}{31}
\newcommand{\RetreatsBackQuietShare}{63}
\newcommand{\PracticalBehind}{250}
\newcommand{\PracticalWithin}{100}
\newcommand{\PracticalCover}{85}
\newcommand{\PracticalMost}{six}
\newcommand{\PracticalDiffersBar}{25}
\newcommand{\PracticalExtraBar}{3}
\newcommand{\PracticalHisPositions}{800}
\newcommand{\PracticalHisKept}{653}
\newcommand{\PracticalHisSOne}{67.1\% {\footnotesize (55.9--76.6)}}
\newcommand{\PracticalHisSZero}{18.0\% {\footnotesize (15.1--21.4)}}
\newcommand{\PracticalHisDiffers}{19.4\% {\footnotesize (16.6--22.7)}}
\newcommand{\PracticalHisExtra}{10.0}
\newcommand{\PracticalAbovePositions}{800}
\newcommand{\PracticalAboveKept}{626}
\newcommand{\PracticalAboveSOne}{56.3\% {\footnotesize (44.8--67.3)}}
\newcommand{\PracticalAboveSZero}{17.9\% {\footnotesize (14.9--21.3)}}
\newcommand{\PracticalAboveDiffers}{20.9\% {\footnotesize (17.9--24.3)}}
\newcommand{\PracticalAboveExtra}{9.1}
\newcommand{\PracticalArchivePositions}{347}
\newcommand{\PracticalArchiveKept}{314}
\newcommand{\PracticalArchiveSOne}{64.3\% {\footnotesize (51.2--75.5)}}
\newcommand{\PracticalArchiveSZero}{22.4\% {\footnotesize (17.7--28.0)}}
\newcommand{\PracticalArchiveDiffers}{22.0\% {\footnotesize (17.8--26.9)}}
\newcommand{\PracticalArchiveExtra}{7.7}
\newcommand{\PracticalCalibration}{%
\setlength{\tabcolsep}{4pt}
\begin{tabular}{llrlrlr}
\toprule
 & \multicolumn{2}{c}{his band} & \multicolumn{2}{c}{the band above} & \multicolumn{2}{c}{his games} \\
\cmidrule(lr){2-3}\cmidrule(lr){4-5}\cmidrule(lr){6-7}
third of Maia's chance & observed & predicted & observed & predicted & observed & predicted \\
\midrule
lowest & 2.8\% {\footnotesize (1.3--5.9)} & 0.0\% & 1.0\% {\footnotesize (0.3--3.5)} & 0.0\% & 1.9\% {\footnotesize (0.5--6.8)} & 0.0\% \\
middle & 0.9\% {\footnotesize (0.3--3.3)} & 0.0\% & 1.0\% {\footnotesize (0.3--3.5)} & 0.0\% & 3.9\% {\footnotesize (1.5--9.6)} & 0.3\% \\
highest & 31.5\% {\footnotesize (25.7--38.0)} & 32.2\% & 32.4\% {\footnotesize (26.4--39.0)} & 32.3\% & 48.1\% {\footnotesize (38.7--57.6)} & 34.5\% \\
\bottomrule
\end{tabular}
}
\newcommand{\PracticalSavedTop}{28.6\% {\footnotesize (23.0--34.9)}}
\newcommand{\PracticalSavedTry}{46.7\% {\footnotesize (30.2--63.9)}}
\newcommand{\PracticalSavedTryCount}{14 of 30}
\newcommand{\PracticalSavedOther}{19.2\% {\footnotesize (15.7--23.3)}}
\newcommand{\PracticalChanceTop}{14.0}
\newcommand{\PracticalChanceTry}{18.0}
\newcommand{\PracticalCheck}{4}
\newcommand{\PracticalPiecesOn}{12}
\newcommand{\PracticalPassed}{2}
\newcommand{\PracticalStalemate}{0}
\newcommand{\LookAbove}{200}
\newcommand{\TripwireWas}{65}
\newcommand{\TripwireWindow}{30}
\newcommand{\GoalsWas}{78\%, 82\% and 85\%}
\newcommand{\ConversionTable}{%
\begin{tabular}{llllrr}
\toprule
 & chess.com & the band it reads (Lichess) & converts & was & now \\
\midrule
stage 1 goal & 1000--1150 & 1450--1599, 200 above & 74.5\% {\footnotesize (73.9--75.2)} & 78\% & 75\% \\
stage 3 goal & 1300--1450 & 1750--1899, 200 above & 75.9\% {\footnotesize (74.7--77.0)} & 82\% & 76\% \\
stage 4 goal & 1450--1600 & 1900--2049, 200 above & 77.9\% {\footnotesize (76.6--79.0)} & 85\% & 78\% \\
tripwire & & 1050--1249, 200 below him & 71.4\% {\footnotesize (70.6--72.2)} & 65\% & 71\% \\
\bottomrule
\end{tabular}
}
\newcommand{\BooksTarget}{70}
\newcommand{\ConvHis}{72.9\% {\footnotesize (72.4--73.5)}}
\newcommand{\ConvAbove}{74.7\% {\footnotesize (74.1--75.3)}}
\newcommand{\ConvTopBand}{1900--2049}
\newcommand{\ConvTop}{77.9\% {\footnotesize (76.6--79.0)}}
\newcommand{\ConvArchive}{74.5\% {\footnotesize (69.6--78.8)}}
\newcommand{\ConvArchiveGames}{337}
\newcommand{\ConvWindow}{30}
\newcommand{\ConvWindowSustained}{18}
\newcommand{\ConvWindowMin}{20}
\newcommand{\PlanMargin}{45}
\newcommand{\PlanPositions}{eleven}
\newcommand{\DeepNodes}{32 million}
\newcommand{\PlanRejected}{5}
\newcommand{\PlanSound}{14}
\newcommand{\PlanRejectedRightPlan}{zero}
\newcommand{\PlanSoundWrongPlan}{zero}
\newcommand{\PlanRejectedRightDeep}{one}
\newcommand{\PlanSoundWrongDeep}{one}
\newcommand{\PlanRejectedTopThree}{one}
\newcommand{\PlanEvalAgreement}{disagrees}
\newcommand{\PlanEvalWhite}{``about even''}
\newcommand{\PlanEvalBlack}{``much worse''}
\newcommand{\LaterAppCommit}{fa07ea1ed8}
\newcommand{\LaterArchiveGames}{647}
\newcommand{\AfterGainMin}{two}
\newcommand{\AfterGainWindow}{three}
\newcommand{\AfterGainLiftBar}{1.1}
\newcommand{\AfterGainClockLow}{a minute}
\newcommand{\AfterGainClockHigh}{three}
\newcommand{\AfterGainHisEvents}{37,151}
\newcommand{\AfterGainHisWindow}{94,543}
\newcommand{\AfterGainHisRate}{8.9\%}
\newcommand{\AfterGainHisExpected}{12.5\%}
\newcommand{\AfterGainHisLift}{0.71 {\footnotesize (0.70--0.73)}}
\newcommand{\AfterGainHisLooseLift}{0.78}
\newcommand{\AfterGainAboveEvents}{28,081}
\newcommand{\AfterGainAboveWindow}{71,410}
\newcommand{\AfterGainAboveRate}{8.4\%}
\newcommand{\AfterGainAboveExpected}{11.6\%}
\newcommand{\AfterGainAboveLift}{0.73 {\footnotesize (0.70--0.75)}}
\newcommand{\AfterGainAboveLooseLift}{0.79}
\newcommand{\AfterGainArchiveEvents}{485}
\newcommand{\AfterGainArchiveWindow}{1,279}
\newcommand{\AfterGainArchiveRate}{8.9\%}
\newcommand{\AfterGainArchiveExpected}{11.3\%}
\newcommand{\AfterGainArchiveLift}{0.79 {\footnotesize (0.63--0.97)}}
\newcommand{\AfterGainArchiveLooseLift}{0.82}
\newcommand{\AfterGainHisMoves}{0.69, 0.70 and 0.75}
\newcommand{\AfterGainHisGains}{0.80 for 2, 0.70 for 3 or 4 and 0.56 for 5 or more}
\newcommand{\AfterGainTable}{%
\setlength{\tabcolsep}{4pt}
\begin{tabular}{lrrllll}
\toprule
 & & & \multicolumn{2}{c}{serious errors} & & \\
\cmidrule(lr){4-5}
players & runs & window moves & in the window & expected & lift & excess (points) \\
\midrule
his band & 37,151 & 94,543 & 8.9\% & 12.5\% & 0.71 {\footnotesize (0.70--0.73)} & $-$3.6 {\footnotesize ($-$3.8 to $-$3.4)} \\
the band above & 28,081 & 71,410 & 8.4\% & 11.6\% & 0.73 {\footnotesize (0.70--0.75)} & $-$3.2 {\footnotesize ($-$3.4 to $-$2.9)} \\
his games & 485 & 1,279 & 8.9\% & 11.3\% & 0.79 {\footnotesize (0.63--0.97)} & $-$2.3 {\footnotesize ($-$4.4 to $-$0.3)} \\
\bottomrule
\end{tabular}
}
\newcommand{\LostTradesBehind}{300}
\newcommand{\LostTradesDown}{two}
\newcommand{\LostTradesDropBar}{3}
\newcommand{\LostTradesMinTrades}{300}
\newcommand{\LostTradesClockEdge}{30}
\newcommand{\LostTradesBothPositions}{11,590}
\newcommand{\LostTradesBothTraded}{33.1\% {\footnotesize (32.3--34.0)}}
\newcommand{\LostTradesBothPlainTrades}{3,008}
\newcommand{\LostTradesBothSavedPlain}{20.3\% {\footnotesize (18.9--21.7)}}
\newcommand{\LostTradesBothSavedKeep}{22.2\% {\footnotesize (21.3--23.2)}}
\newcommand{\LostTradesBothDiff}{$-$1.7 {\footnotesize ($-$3.4 to 0.0)}}
\newcommand{\LostTradesHisPositions}{6,370}
\newcommand{\LostTradesHisTraded}{33.3\% {\footnotesize (32.1--34.5)}}
\newcommand{\LostTradesHisPlainTrades}{1,657}
\newcommand{\LostTradesHisSavedPlain}{21.5\% {\footnotesize (19.6--23.5)}}
\newcommand{\LostTradesHisSavedKeep}{23.6\% {\footnotesize (22.3--25.0)}}
\newcommand{\LostTradesHisDiff}{$-$1.9 {\footnotesize ($-$4.4 to 0.5)}}
\newcommand{\LostTradesAbovePositions}{5,220}
\newcommand{\LostTradesAboveTraded}{32.9\% {\footnotesize (31.6--34.1)}}
\newcommand{\LostTradesAbovePlainTrades}{1,351}
\newcommand{\LostTradesAboveSavedPlain}{18.7\% {\footnotesize (16.7--20.9)}}
\newcommand{\LostTradesAboveSavedKeep}{20.5\% {\footnotesize (19.2--22.0)}}
\newcommand{\LostTradesAboveDiff}{$-$1.4 {\footnotesize ($-$4.0 to 1.1)}}
\newcommand{\LostTradesArchivePositions}{75}
\newcommand{\LostTradesArchiveTraded}{37.3\% {\footnotesize (27.3--48.6)}}
\newcommand{\LostTradesArchivePlainTrades}{23}
\newcommand{\LostTradesArchiveSavedPlain}{21.7\% {\footnotesize (9.7--41.9)}}
\newcommand{\LostTradesArchiveSavedKeep}{28.6\% {\footnotesize (17.2--43.6)}}
\newcommand{\LostTradesArchiveDiff}{2.5 {\footnotesize ($-$27.8 to 29.4)}}
\newcommand{\LostTradesBothDrop}{1.7}
\newcommand{\LostTradesAfterDiff}{$-$3.4 {\footnotesize ($-$5.2 to $-$1.7)}}
\newcommand{\LostTradesClockDiff}{$-$2.3 {\footnotesize ($-$4.1 to $-$0.5)}}
\newcommand{\LostTradesLossPlain}{1.4}
\newcommand{\LostTradesLossKeep}{2.6}
\newcommand{\LostTradesCheck}{786}
\newcommand{\LostTradesCheckDiff}{$-$1.7 {\footnotesize ($-$4.7 to 1.4)}}
\newcommand{\LostTradesPassed}{42}
\newcommand{\LostTradesPawn}{90}
\newcommand{\LostTradesPawnDiff}{$-$8.4 {\footnotesize ($-$15.1 to $-$0.9)}}
\newcommand{\LostTradesOther}{843}
\newcommand{\LostTradesOtherDiff}{0.2 {\footnotesize ($-$3.0 to 3.5)}}
\newcommand{\LostTradesMoveLoses}{six}
\newcommand{\LostTradesAll}{3,836}
\newcommand{\LostTradesTable}{%
\setlength{\tabcolsep}{4pt}
\begin{tabular}{lrrrrll}
\toprule
 & & & \multicolumn{2}{c}{saved after} & \multicolumn{2}{c}{plain trade $-$ keeping, in points} \\
\cmidrule(lr){4-5}\cmidrule(lr){6-7}
players & positions & traded & a plain trade & keeping the pieces & alike before & alike after \\
\midrule
both bands & 11,590 & 33\% & 20.3\% & 22.2\% & $-$1.7 {\footnotesize ($-$3.4 to 0.0)} & $-$3.4 {\footnotesize ($-$5.2 to $-$1.7)} \\
his band & 6,370 & 33\% & 21.5\% & 23.6\% & $-$1.9 {\footnotesize ($-$4.4 to 0.5)} & $-$3.5 {\footnotesize ($-$6.0 to $-$1.1)} \\
the band above & 5,220 & 33\% & 18.7\% & 20.5\% & $-$1.4 {\footnotesize ($-$4.0 to 1.1)} & $-$3.4 {\footnotesize ($-$5.9 to $-$0.9)} \\
his games & 75 & 37\% & 21.7\% & 28.6\% & 2.5 {\footnotesize ($-$27.8 to 29.4)} & 5.1 {\footnotesize ($-$21.3 to 31.7)} \\
\bottomrule
\end{tabular}
}
\newcommand{\VerdictTable}{%
\begin{tabular}{lp{0.42\textwidth}p{0.42\textwidth}}
\toprule
test & what the rule found & what the app took \\
\midrule
1a & Traps are a real share of the failures, and misses dominate & A ``take it or not'' drill (chess PR 109); the punish goal keeps its hanging-piece practice: no other kind reaches one miss in ten \\
1b & Four of eight signs warn; past the book, two & Practice games warn on their reply (chess PR 111) \\
1c & No case tightens the praise & Nothing: the praise stays \\
1d & Backward moves are missed less often, not more, and are not harder for Maia & Nothing \\
1e & Maia is not a fair judge here; the practical try does not save more & Nothing: the review names the top move \\
1f & All four conversion numbers move & The goals and the tripwire (chess PR 108) \\
1g & No disagreement holds & Nothing: the margin stays \\
6a & Serious errors are rarer, not more common, right after winning material & The scan drops its clause on winning material \\
6b & An even trade when lost costs fewer saved games than the bar & Nothing \\
\bottomrule
\end{tabular}
}


\title{Does the books' advice hold at his level?\\[0.3em]
\large Nine rules from two chess books, tested against Stockfish and Maia-2 on Lichess 10-minute
games and his own}
\author{Chess Coach research lab}
\date{}

\begin{document}
\maketitle

\begin{abstract}
Two chess books for club players answered twenty questions the coaching app raised. Where an answer
suggests a rule the app could follow or a number it could aim at, this paper tests it before the app
acts on it, each under a rule written down and committed before its run, on \GamesRead{} rated
10-minute Lichess games and his \ArchiveGames. After the opponent's serious error, \OneAHisPunishedShare\% of
replies at his level punish it; \OneAHisTrapShare\% of the failures take material that is a trap, enough
for a drill that mixes gifts with traps, and the rest are misses whose best reply is rarely a fork, a
mate or a pin, so the punish goal's practice stays on hanging pieces. Four of eight warning signs of danger
to one's own king clear the bar at his level, but past the first \BookPlies{} plies, where a warning
in a game can speak, only two hold: the enemy queen on the king's exits and a king with no safe
square (lifts \KingLateQueenOnExits{} and \KingLateNoSafeSquares); the rook lift and two more attackers
than defenders mostly mark a later stage of the game. The coach's ``good trade'' praise is faulted
on \TradeHisFaulted{} of praised trades at his level, and no case the books warn about makes it
worse. Backward moves are not the moves his level misses most: a quiet backward best move is missed
less often than a quiet forward one, and Maia-2 finds it more likely. In lost positions, Maia-2 is not a fair judge of which move sets the most
problems, so the review keeps naming the engine's top move. The books' \BooksTarget\% conversion
target is already passed by players at his level; the app's three conversion goals were above what
players \LookAbove{} points above each stage convert, and move down to those rates. On the books'
twelve planning positions, the guess-the-plan drill's \PlanMargin{} cp rule agrees with every
verdict at its own budget. Two later tests put the books' warnings about winning and losing positions
to the same kind of rule. Serious errors are rarer, not more common, in the three moves after a
player wins material than in positions alike, so the scan stops singling out that moment. When two
pawns down and losing, an even trade of pieces saves \LostTradesBothDrop{} fewer games in 100 than
keeping the pieces, under the bar. Compared at the same eval after the move it saves
\LostTradesAugAfterDrop{} fewer, past the bar on a second month of games, so the live coach now says
so where the engine agrees.
\end{abstract}

\section{The question}

Levy Rozman's \emph{How to Win at Chess} (2023) and \emph{How to Win at Chess: Next Level} (2026)
answered twenty questions about the app's coaching. Most answers are qualitative. Seven suggest a
rule or a number, and each became a test (Table~\ref{tab:verdicts}):
\begin{enumerate}[nosep,label=\textbf{1\alph*}]
\item After the opponent's serious error, how often do players at his level punish it, miss it, or
take material that is a trap?
\item Which warning signs of danger to one's own king does the engine bear out?
\item How often does the coach's ``good trade'' praise land on a trade the books would not praise?
\item Are retreats, and backward knight moves above all, the best moves his level misses most, and
the hardest for Maia to find?
\item When every move loses, does the move whose winning reply is hardest to find save more games
than the engine's top move?
\item What share of decisive advantages do players at his level and \LookAbove{} points above
convert, against the app's goals?
\item Does the guess-the-plan drill's \PlanMargin{} cp rule agree with the books' verdicts on their
twelve planning positions?
\end{enumerate}
A second look at the answers gave two more:
\begin{enumerate}[nosep,label=\textbf{6\alph*}]
\item Are serious errors more common in the three moves after a player wins material than in
positions alike?
\item When down material and losing, does an even trade of pieces lower the chance of saving the
game?
\end{enumerate}
No model judges anything here. Every verdict comes from Stockfish, Maia-2 or the app's own board
code, imported through the lab's bridge at app commit \texttt{\AppCommit} (6a and 6b, run later, at
\texttt{\LaterAppCommit}), and each test's rule, its bar and what the app does when it passes were
committed before its run (the project's \texttt{README.md}).

\begin{table}[h]
\centering
\small\VerdictTable
\caption{The nine tests, what each rule decided, and what the app took.}
\label{tab:verdicts}
\end{table}

\section{The data}

The games are the \GamesRead{} rated 10-minute Lichess games of \Month{} (CC0) that the
think-time study keeps, with Lichess's evaluation on every move and the clock. Each gives two
player-games, one for each player, read the app's way. A player-game's band is the player's own
Lichess rating. His chess.com rapid rating comes to the middle of his band, \HisLo--\HisHi, by the
app's conversion; the band above is \AboveLo--\AboveHi. His band holds \HisPlayerGames{}
player-games and the band above \AbovePlayerGames. His games are his \ArchiveGames{} 10-minute games
in the app's archive at that commit; only aggregates are committed. The engine is the app's Node
engine service (Stockfish 18 lite, one thread a search), at the app's own budgets, and Maia-2 comes
through the same service. The books' positions (test 1g) are not in this repository: the test reads
them at run time and commits counts only. Tests 6a and 6b need no engine: they read the games' own
evaluations and the app's board code, and at their commit his archive holds \LaterArchiveGames{}
10-minute games. The recheck of 6b reads a second month, August 2026, sampled and built the same
way: \LostTradesAugGames{} player-games in the two bands.

Shares carry 95\% Wilson intervals.

\section{1a. The move after their mistake}

The books warn that a piece offered for free may be a trap, and that the usual process applies after
the opponent's mistake: check what changed, look at the forcing moves, test the reply. The app's
``Punish their mistakes'' goal drills hanging-piece puzzles.

\textbf{The rule.} The event is a player's move right after the opponent's serious error (a mistake
or a blunder by the review reader). The reply \emph{punished} it when it kept at least half of the
gift, the app's own test; it took a \emph{trap} when it did not and was a capture the reader judged a
mistake or a blunder; otherwise it \emph{missed}. Traps are a real share when at least one failure
in ten is a trap in his band, or in his games with the interval's lower end there or above; then the
app gets a ``take it or not'' drill and the review tags a poisoned capture. Misses dominate when at
least two failures in three are misses in his band; then the goal's practice adds the punishments it
lacks, by what the misses missed.

\textbf{The split.} Both hold (Table~\ref{tab:onea}). At his level \OneAHisPunished\% of replies
punish the error. Of the failures, \OneAHisTrapOfFailures{} are a trap taken and
\OneAHisMissedOfFailures{} a miss; his own games split the same way (\OneAArchiveTrapOfFailures{}
traps). On the punish goal's own population, gifts of \GiftFloor{} win\% or more, traps are a larger share
of the failures (\OneAHisGiftTrapOfFailures{} at his level). After a punished reply, the player's next
move is a serious error \OneAHisNextSerious{} of the time, against \OneAHisHighOtherSerious{} on his
other moves at \HighWp{} win\% or more: the books' ``success is not permission to stop looking''.

\begin{table}[h]
\centering
\scriptsize\OneATable
\caption{The reply to the opponent's serious error.}
\label{tab:onea}
\end{table}

\textbf{What the misses missed.} The engine (\BestNodes{} nodes) found the best reply in a seeded
sample of \OneAHisMissesSearched{} misses in each band and in all \OneAArchiveMissesSearched{} of
his, and the app's labeller named it (Table~\ref{tab:oneakinds}). No kind of punishment beyond the
hanging piece names one miss in ten in his band: a mate \OneAHisKindMate, a fork \OneAHisKindFork,
a pin, skewer or discovered attack \OneAHisKindPinSkewerDiscovered. The punish goal keeps its
hanging-piece practice. The free piece itself is the best reply in \OneAHisKindHangingPiece{} of the
misses and saving a loose piece of the player's own in \OneAHisKindDefensiveMove; the labeller does
not name the rest (\OneAHisKindUnnamed).

\begin{table}[h]
\centering
\scriptsize\OneAKindsTable
\caption{The engine's best reply in the position the opponent's serious error left, where the player
missed the punishment, as the app's labeller names it.}
\label{tab:oneakinds}
\end{table}

\textbf{Offered material.} In his band, \OneAHisOfferedSearched{} sampled replies came right after
the opponent's move left one of their knights, bishops, rooks or queens newly hanging. Taking it with
the cheapest attacker was a trap, a mistake or a blunder by the app's judge, in
\OneAHisOfferedTraps{} of them, most often for a knight, \OneAHisOfferedTrapKnight, and least
often for a queen, \OneAHisOfferedTrapQueen. The players took \OneAHisTookTrap{} of the traps and left
\OneAHisLeftGift{} of the real gifts, so leaving a gift was the commoner failure. In his games
\OneAArchiveOfferedTraps{} of \OneAArchiveOfferedSearched{} offers were traps.

\section{1b. Danger to his own king}\label{sec:king}

The books list warning signs that come before a mate is on the board: a king with no safe squares,
an enemy queen covering its exits, a bishop aimed into the shelter, an approaching knight, a rook
lift, a hook pawn, a pawn storm when the kings stand on opposite wings, and at least two more
effective attackers than defenders near the king (\emph{plus two}).

\textbf{The rule.} The positions are one position before a player's move from each of \KingHisPositions{}
player-games in each band, past the book, not in check and not an endgame by the app's phase divider,
and up to ten from each of his games. A position is \emph{dangerous} when the engine finds a forced
mate for the opponent with the opponent to move (\CheckNodes{} nodes), or the game's own evals put
the player \SwingWp{} win\% or more lower within \SwingPlies{} plies. Each sign is board code over the app's board kernel,
over the king's zone (its square, the eight around it and the three squares two ranks in front). A
sign becomes a warning when, in his band, it shows in at least \KingShowsBar\% of positions, lifts the
danger rate at least \KingLiftBar{} times, and its precision's interval starts at \KingFloorBar{}
times the base rate or more.

\textbf{The result.} \KingClearedCap{} signs clear the bar (Table~\ref{tab:king}). The queen on the exits
lifts the danger rate most. The opposite-side storm lifts it \KingStormLift{} times, but its interval
starts at \KingStormLow\%, under the floor of \KingFloor\%. No safe squares foretells mate threats
(lift \KingNoSafeMateLift) and not swings (\KingNoSafeSwingLift).

\begin{table}[h]
\centering
\scriptsize\KingTable
\caption{The eight signs in his band's positions (base rate: mate threats \KingHisMateBase\%, swings
\KingHisSwingBase\%), with each sign's lift in the band above and in his games.}
\label{tab:king}
\end{table}

\textbf{Where a warning can speak.} A step taken after the result. Inside the first \BookPlies{}
plies the signs say nothing: any sign was dangerous \KingEarlyAny{} of the time against a base of
\KingEarlyBase\% (\KingEarlyPositions{} positions), since the uncastled king, boxed in by its own
pieces, shows ``no safe squares'' from the first move. A practice game from the start can only tell
the book is over by waiting out those plies, so the same bar was applied again to the
\KingLatePositions{} positions after them (Table~\ref{tab:kinglate}). Danger rises with the ply, from
\KingBinFirst\% in plies \KingBinFirstPlies{} to \KingBinLast\% from ply \KingBinLastFrom, and the rook lift and plus two mostly
marked that rise: past the book their lifts are \KingLateRookLift{} and \KingLatePlusTwo, and against
the danger each ply bin predicts they lift it little. Only the queen on the exits and no safe squares
clear the bar there, and they keep a lift of \KingLiftBar{} or more against the ply alone in every group. This
step could take signs away, never add one.

\begin{table}[h]
\centering
\scriptsize\KingLateTable
\caption{The four signs that cleared the bar, on his band's positions after the first \BookPlies{}
plies. ``Against the ply alone'' is the danger observed over the danger expected from the base rate of
each ply bin (\KingBins).}
\label{tab:kinglate}
\end{table}

In his games, after the first \BookPlies{} plies, a warning on the two signs, said when a sign appears
that the position before his last move did not show, would speak \KingNotesPerGame{} times a game
(none in \KingNotesSilent\% of games; \KingNotesHisBand{} in his band's). Danger follows it
\KingOnsetDanger{} of the time, against \KingNoneDanger{} where no sign shows. On all four signs it
would speak \KingNotesAllPerGame{} times a game, followed by danger \KingAllOnsetDanger{} of the time.

\section{1c. The coach's ``good trade'' praise}

When the player is ahead, the live coach praises an even trade: ``Good trade: knight for knight keeps
you up a piece, and with fewer pieces left that lead counts for more.'' The books qualify the advice:
do not give a more active piece for a less active one, nor a piece one's attack needs, nor the queens
while attacking, and judge the ending the trade leads to.

\textbf{The rule.} Every move where the coach's trade rule fires in his games, and a sample of
\TradeHisTrades{} in each band, is read as the praise gate reads it, at the coach's own budgets; the
praised trades are the population. A case tightens the gate when it holds for at least one praised
trade in twenty (\TradeShareBar\%) in his band or in his games, and the game's own eval faults its
praised trades (a loss of \TradeFaultBar{} win\% or more) at least as often as the others. A piece is
\emph{more active} when it has \TradeMobilityGap{} more legal moves or more. A trade leads to a \emph{worse
ending} when the player's next move is in the endgame and the game's own eval faults the trade, so
every trade in that case is faulted and only its share can decide.

\textbf{The result.} No case tightens the praise (Table~\ref{tab:trades}). The coach praises
\TradeHisPraised\% of the trades its rule finds at his level, and the game's eval faults
\TradeHisFaulted{} of them (\TradeHisFaultedCount); in his games \TradeArchiveFaulted{}
(\TradeArchiveFaultedCount). Giving the more active piece is common, \TradeActiveShare\% of the
praised trades, but those are faulted less often than the rest (\TradeActiveFaulted\% against
\TradeActiveOthers\%). Queens off while attacking is faulted more (\TradeQueensFaulted\%), on
\TradeQueensShare\% of praised trades, under the one in twenty the rule asks. The review's
missed-trade line fires on \TradeLineFires{} of his \TradeLineFaultedAhead{} faulted moves while
ahead; the best move's trade gives the more active piece on \TradeLineActive, trades an attacker on
\TradeLineAttacker{} and takes the queens off while attacking on \TradeLineQueens.

\begin{table}[h]
\centering
\scriptsize\TradeTable
\caption{The books' cases among the coach's praised trades: each case's share, how often the game's
own eval faults the praised trades in it, and the other praised trades.}
\label{tab:trades}
\end{table}

\section{1d. Backward moves}

The books call backward knight moves among the hardest moves to notice, and ask the solver to check
whether a retreat helps one or two moves later.

\textbf{The rule.} The positions are the position before a player's move, past the book, with at
least two legal moves: up to ten from each of \RetreatsDrawGames{} player-games in each band, and
\RetreatsDrawArchive{} from his games. Only \emph{must-find} positions count, where the engine
(\CheckNodes{} nodes, two lines) finds a best move at least \RetreatsMustFind{} win\% better than the
second, so that missing it is at least a mistake by the app's judge. The player \emph{found} it when
he played the engine's best move. A best move is \emph{backward} when it ends on a rank nearer the
mover's own back rank, \emph{forward} when nearer the opponent's, and \emph{sideways} otherwise;
castling is left out. The tests compare \emph{quiet} best moves (no capture, no promotion), since a
capture is seen as a capture whichever way it goes. Backward moves are \emph{missed more often} when,
in his band, the miss rate of backward quiet best moves exceeds that of forward ones with the
difference's interval (Newcombe's) above zero. They are \emph{harder for Maia at the same value} when
Maia-2's probability for the best move, at the player's band against the opponent's, averages lower
for backward than for forward quiet best moves within each of three bins of the best move's margin
over the second (\RetreatsBins{} win\%), and the difference, weighted by the backward moves in each
bin, has a bootstrap interval below zero. When both hold, the scan adds ``look at backward moves
too'' and the review names a missed retreat.

\textbf{The result.} Neither half holds, and both point the other way (Table~\ref{tab:retreats}). In
his band a quiet backward best move was missed \RetreatsHisBackward{} of the time, against
\RetreatsHisForward{} for a quiet forward one, a difference of \RetreatsHisDiff{} points. The band
above is the same, and in his games the two are missed about equally. Backward knight moves were
missed \RetreatsHisKnightBack{} of the time against \RetreatsHisKnightFwd{} for forward knight moves,
inside the noise. Maia-2 gives a backward best move more probability than a forward one at the same
margin, not less (Table~\ref{tab:retreatsmaia}), in the two bins that hold
\RetreatsFirstBinsBackward{} of the \RetreatsHisQuietBackward{} backward moves; only in the last,
with \RetreatsLastBinBackward{} of them, does it give less. With captures counted, backward and
forward best moves are missed about equally (\RetreatsHisAllBackward{} against
\RetreatsHisAllForward), since only \RetreatsFwdQuietShare\% of forward best moves are quiet against
\RetreatsBackQuietShare\% of backward ones, and a capture is found more often. The move his level
misses most is the quiet move forward. The scan and the review stay as they are.

\begin{table}[h]
\centering
\scriptsize\RetreatsTable
\caption{How often the player missed a must-find best move, by the direction of a quiet best move.
Differences are in points, with Newcombe's interval.}
\label{tab:retreats}
\end{table}

\begin{table}[h]
\centering
\scriptsize\RetreatsMaiaTable
\caption{Maia-2's mean probability for the best move in his band, by the best move's margin over the
second, with the bootstrap interval of the weighted difference.}
\label{tab:retreatsmaia}
\end{table}

\section{1e. Practical tries when every move loses}

The books say that in a lost position the defender should make the opponent's task hard. Where every
move leaves the player behind and he picked a worse one, the review says ``You are behind either
way'' and names the engine's top move.

\textbf{The rule.} One position from each player-game where the game's own eval puts the player
\PracticalBehind{} cp or more behind (\PracticalHisPositions{} in each band, every one in his games),
kept when the engine's top move still loses. The candidates are the engine's lines within
\PracticalWithin{} cp of the top move. For each, Maia-2 at the opponent's level gives the opponent's
likeliest replies (up to \PracticalMost, until they cover \PracticalCover\%), and the engine scores
them; a reply \emph{throws the win away} when it leaves the player no longer behind. A candidate's
\emph{chance} is Maia's probability on those replies, and the \emph{practical try} is the candidate
with the highest chance. The practical try saves more games when Maia is a fair instrument (the
replies that throw the win away rise from the lowest third of Maia's chance to the highest, and the
predicted chance lies inside the observed share in at least two thirds), it differs from the top
move in at least \PracticalDiffersBar\% of his band's positions, and its extra saves average at
least \PracticalExtraBar{} games per 100 there.

\textbf{The result.} It does not. Maia is not a fair instrument here (Table~\ref{tab:practical}): in
two thirds of the positions it gives no chance at all, yet some replies there do throw the win away;
in the top third it is right. The practical try differs from the top move in \PracticalHisDiffers{}
of his band's \PracticalHisKept{} positions, under the bar. Where it differs, its extra saves would be
\PracticalHisExtra{} games per 100, from a mean chance of \PracticalChanceTry\% against
\PracticalChanceTop\% for the top move. The game's own outcomes agree in direction: players who
played the practical try were saved \PracticalSavedTry{} of the time (\PracticalSavedTryCount), the
top move \PracticalSavedTop, any other move \PracticalSavedOther. Where it differs, the practical try
gives a check in \PracticalCheck\% of positions, keeps pieces on in \PracticalPiecesOn\%, pushes a
passed pawn in \PracticalPassed\% and stalemates in \PracticalStalemate\%.

\begin{table}[h]
\centering
\scriptsize\PracticalCalibration
\caption{Maia's chance that the opponent's reply throws the win away, by third of the predicted
chance, against the share of actual replies that did.}
\label{tab:practical}
\end{table}

\section{1f. The conversion target}

The books' informal target is converting a decisive advantage at least \BooksTarget\% of the time.
The app measures conversion as \texttt{conversion.sustained}: of the games where the player stood at
+3 or better for three plies in a row, the share won. Its tripwire fires under \TripwireWas\% over
the trailing \TripwireWindow{} games, and its three stage goals ask for \GoalsWas.

\textbf{The rule.} A stage goal stays when it lies inside the interval of what players \LookAbove{}
points above that stage's band convert, and otherwise moves to that rate; the tripwire should fire
when he converts no better than players \LookAbove{} points below him.

\textbf{The result.} All four move (Table~\ref{tab:conversion}). Each goal was above what players
\LookAbove{} points above its stage convert, and the stage 4 goal asked for more than Lichess players
at \ConvTopBand{} convert, \ConvTop. Players at his level already pass the books' \BooksTarget\%,
converting \ConvHis, and the band above converts \ConvAbove. In his \ConvArchiveGames{} games with a sustained +3 he won \ConvArchive.
His last \ConvWindow{} games at that commit hold \ConvWindowSustained{} with a sustained +3,
under the metric's minimum of \ConvWindowMin, so the tripwire stays quiet until it has the games.

\begin{table}[h]
\centering
\scriptsize\ConversionTable
\caption{The conversion goals and the tripwire, the band each is read from, and what it moves to.}
\label{tab:conversion}
\end{table}

\section{1g. The twelve planning positions}

The guess-the-plan drill shows a quiet position with the engine's three best moves and counts any
option within \PlanMargin{} cp of the best as right. Next Level's last chapter gives twelve positions
with the author's preferred, sound and rejected moves: \PlanRejected{} rejected and \PlanSound{}
preferred or sound over \PlanPositions{} positions, and one that asks for an evaluation.

\textbf{The rule.} Each position is searched at the drill's own budget (\PlanNodes{} nodes) and at
\DeepNodes{} nodes. A
disagreement is a rejected move the rule counts as right, or a preferred or sound move it counts as
wrong, and it holds when both budgets give it.

\textbf{The result.} None holds. At the drill's budget the rule counts \PlanRejectedRightPlan{}
rejected moves as right and \PlanSoundWrongPlan{} preferred or sound moves as wrong; the deep check
alone gives \PlanRejectedRightDeep{} and \PlanSoundWrongDeep. Of the rejected moves,
\PlanRejectedTopThree{} is among the drill's three options at all. The evaluation position, which the books call equal
without saying who is to move, the engine reads at \DeepNodes{} nodes as \PlanEvalWhite{} for White
with White to move and \PlanEvalBlack{} for White with Black to move, so it \PlanEvalAgreement. The
margin stays.

\section{6a. Errors right after winning material}

The books warn that a tactical success is no licence to stop looking at the opponent: the author's own
win of material against a strong player was followed by a forcing reply he overlooked. The app's
scan ends: ``Before you move: what is wrong with your move, and what is their best answer? Ask it
hardest right after you win material.''

\textbf{The rule.} A \emph{capture run} is a stretch of consecutive plies that each capture or
promote. A run \emph{wins material} for a player when he captures or promotes in it and the material
count (pawn 1, knight and bishop 3, rook 5, queen 9) stands \AfterGainMin{} or more higher for him
after it than before: a piece for a pawn, the exchange, a free piece or a new queen kept, but not a
pawn or an even trade. The \emph{window} is his next \AfterGainWindow{} moves after his last
capture in the run, and the controls are his other moves, all past the book with a choice to make.
A window move's expected chance of a serious error (the review reader's mistake or blunder) is the
controls' rate in its stratum: his win\% before the move in ten-point bins, tension (in check, right
after a capture, or a piece hanging), either of the two signs of danger to his king that Play warns
on, the phase, and his clock (under \AfterGainClockLow, one to \AfterGainClockHigh{} minutes,
\AfterGainClockHigh{} or more). It holds when, in his band, the window's serious errors exceed that
expectation with a bootstrap interval above zero and their ratio, the \emph{lift}, is at least
\AfterGainLiftBar. Then the routines' brief would get a reading for it; otherwise the scan drops its
last sentence.

\textbf{The result.} It does not hold, and it points the other way (Table~\ref{tab:aftergain}). In
his band \AfterGainHisEvents{} runs won material, and \AfterGainHisRate{} of the
\AfterGainHisWindow{} window moves were serious errors, against \AfterGainHisExpected{} expected in
positions alike: a lift of \AfterGainHisLift. The band above and his games read the same way. Each
of the three window moves is lower on its own (lifts \AfterGainHisMoves), and the more material
won, the lower the lift (\AfterGainHisGains{} pawns won). On the eval and the two dangers alone the
lift is \AfterGainHisLooseLift: the phase and the clock of the window moves, where errors are more
common, widen the gap rather than close it. The moment after winning material is not where errors
gather at his level, so the scan's last step drops ``Ask it hardest right after you win material''
and keeps its question for every move.

\begin{table}[h]
\centering
\scriptsize\AfterGainTable
\caption{Serious errors in the three moves after a run that wins material, against their expected
number at the rate of other moves in positions alike. Lift and excess with bootstrap intervals over
the player-games.}
\label{tab:aftergain}
\end{table}

\section{6b. Even trades when lost}

The books' chapter on defending discourages even trades when down material, since fewer pieces leave
fewer chances to resist, and allows the trade that improves the position, makes counterplay or makes
a passed pawn. Before this test the app said nothing about trades when behind; it praised a trade
only when ahead (1c).

\textbf{The rule.} In each player-game, the first position where the player, past the book, not in
check and after an opponent move that captured nothing, stands \LostTradesBehind{} cp or more behind
by the game's own eval and \LostTradesDown{} pawns or more down in material, with an even trade of
pieces on offer: a capture of a knight, bishop, rook or queen whose exchange nets nothing, by the
app's own trade reader. A \emph{plain trade} takes it, without check and without making a passed
pawn; \emph{keeping the pieces} is a move that captures nothing; trades with check, trades that make
a passed pawn, even pawn trades and other captures are counted apart. A game is \emph{saved} when the
player drew or won. Within each stratum (his win\% before the move in five-point bins, the material
down, his rating in 100-point bins, and whether the opponent is rated higher), the save rate after a
plain trade minus after keeping the pieces, weighted by the plain trades. It holds when, over both
bands together, that difference's bootstrap interval lies below zero and the drop is at least
\LostTradesDropBar{} saved games in 100. Then Play's coach and the review would say so where the
engine agrees.

\textbf{The result.} It does not hold (Table~\ref{tab:losttrades}). In \LostTradesBothPositions{}
positions the players traded \LostTradesBothTraded{} of the time. A plain trade saved
\LostTradesBothSavedPlain{} of its games and keeping the pieces \LostTradesBothSavedKeep; in positions
alike the difference is \LostTradesBothDiff{} points, under the bar, with an interval that reaches
zero. The trades themselves gave away less, \LostTradesLossPlain{} win\% against
\LostTradesLossKeep{} for the moves that keep the pieces, so at the same eval after the move the
trade costs more: \LostTradesAfterDiff{} points, reported without a rule. With the clocks as a
further stratum it is \LostTradesClockDiff. Trades with check save about as often as keeping the
pieces, at \LostTradesCheckDiff{} points, the \LostTradesPawn{} even pawn trades saved fewer, at
\LostTradesPawnDiff{} points, and the \LostTradesPassed{} trades that made a passed pawn are too few
to read. His games hold \LostTradesArchivePositions{} such positions, too few to read. By the rule,
nothing changes.

\textbf{After the result.} The measure at the same eval after the move was not the rule's, so it was
tested again by a rule fixed before that run, on games the study had not read: the Lichess games of
August 2026, streamed and built as September's at the same app commit, with the same positions,
moves and outcome and the strata of that measure (his win\% after the move in five-point bins, the
material down, his rating, and whether the opponent is rated higher). It holds when, over both
bands, the interval lies below zero and the drop is at least \LostTradesDropBar{} saved games in
100, with at least \LostTradesMinTrades{} plain trades. It holds. Over \LostTradesAugPlain{} plain
trades the difference is \LostTradesAugAfterDiff{} points, \LostTradesAugHisAfter{} in his band and
\LostTradesAugAboveAfter{} in the band above. The rule's own measure on those games is
\LostTradesAugBeforeDiff{} points, below zero this time and still under the bar. So the live coach
now speaks when he is behind and takes an even trade of pieces that the engine scores no better than
a move that keeps them, and the finding goes to the review's explanations.

\begin{table}[h]
\centering
\scriptsize\LostTradesTable
\caption{At the first position where a player two pawns or more down and losing could trade pieces
evenly: how often he traded, how often the game was saved, and the difference in positions alike by
the eval before the move (the rule) and after it, in September's games and in August's, which
rechecked the measure after the move.}
\label{tab:losttrades}
\end{table}

\section{Limits}

The bands are other players, in games against their own level: their numbers say where the books'
rules hold at his level, not what he would gain by following them. Danger in 1b is read from Lichess's
evaluations six plies on and a mate search at a modest budget, so it misses slow attacks. The step
after 1b's result applied a bar fixed before the run to a population chosen after it; it only took
signs away. In 1a the labeller names one theme per best reply, has no trapped piece, and calls what
it cannot name by the phase. Maia-2's probabilities in 1e cover the likeliest
replies only, which is where it fails the test. In 6a the strata match what the app can read; a move
right after winning material may still be easier in ways they do not hold. In 6b the trade reader
judges a capture by its square, so \LostTradesMoveLoses{} of the \LostTradesAll{} trades lost
material as played. Its measure after the move was chosen after September's result; August's games
tested it fresh, but both months come from one site and one time control, and the save rate counts
what players did, not what keeping the pieces would have done for the same player. His Lichess level is the app's conversion of his chess.com rating,
and his games come from the app's reviews, with another engine and depth.

\end{document}
